\chapter{Summary of Notation and Rules}
\label{app:summary}

\section*{Syntax}

\begin{center}
\begin{tabular}{@{}ll@{}}
\toprule
Symbol & Meaning\\
\midrule
$\cir{o\; t_1\cdots t_n}$ & circle with operator $o$ and operands $t_1,\dots,t_n$\\
$\Num,\ \Color,\ \Image,\ \Sym$ & sorts of the core fragment; $\Nat$ is introduced in Chapter~\ref{ch:sorts}\\
$\delta_o$ & the partial function giving the meaning of operator $o$\\
$\dom(\delta_o)$ & the admissible operand tuples of $o$\\
$\size(t)$ & number of circles in $t$\\
$\depth(t)$ & greatest nesting of circles in $t$\\
$t|_p$ & subterm of $t$ at position $p$\\
$t[p\leftarrow s]$ & $t$ with the subterm at $p$ replaced by $s$\\
$\sigma(t)$ & substitution $\sigma$ applied to $t$\\
$t\pop t'$ & $t$ pops to $t'$ in one step\\
$t\popstar t'$ & $t$ pops to $t'$ in zero or more steps\\
$\nf(t)$ & the normal form of $t$\\
$s\equiv t$ & $s$ and $t$ have the same normal form\\
\bottomrule
\end{tabular}
\end{center}

\section*{Reduction}

\begin{align*}
\cir{o\; v_1\cdots v_n} &\pop \delta_o(v_1,\dots,v_n) &&\text{if }(v_1,\dots,v_n)\in\dom(\delta_o) &&\text{(Pop)}\\
\cir{o\; \dots t_i\dots} &\pop \cir{o\; \dots t_i'\dots} &&\text{if } t_i\pop t_i' &&\text{(Arg)}
\end{align*}

\section*{Operators used in the text}

\begin{center}
\begin{tabular}{@{}lll@{}}
\toprule
Operator & Signature & Admissible when\\
\midrule
$\op{add},\ \op{sub},\ \op{mul}$ & $\Num\times\Num\to\Num$ & always\\
$\op{pow}$ & $\Num\times\Num\to\Num$ & exponent is non-negative\\
$\op{rotate}$ & $\Image\times\Num\to\Image$ & angle is a multiple of $90$\\
$\op{recolor}$ & $\Image\times\Color\to\Image$ & always\\
$\op{beside}$ & $\Image\times\Image\to\Image$ & images have equal height\\
\bottomrule
\end{tabular}
\end{center}

\section*{Main results}

\begin{center}
\begin{tabular}{@{}lp{0.68\textwidth}@{}}
\toprule
Result & Statement\\
\midrule
Proposition~\ref{prop:size-drop} & Each pop lowers the size by exactly one.\\
Theorem~\ref{thm:termination} & Every sequence of pops is finite, of length at most the size.\\
Lemma~\ref{lem:diamond} & Two different pops from the same term join in one step each.\\
Theorem~\ref{thm:unique} & Every term has a unique normal form.\\
Corollary~\ref{cor:order} & All maximal pop sequences have the same length.\\
Proposition~\ref{prop:parallel} & Parallel popping needs exactly as many rounds as the depth.\\
Theorem~\ref{thm:preservation} & Popping preserves sorts.\\
Theorem~\ref{thm:progress} & A well-sorted blocked term has an inadmissible innermost circle.\\
Theorem~\ref{thm:staged} & Evaluating with symbols first, then substituting, agrees with substituting first.\\
Theorem~\ref{thm:rewrite-conf} & The arithmetic simplification rules are terminating and confluent.\\
\bottomrule
\end{tabular}
\end{center}
